% Extract: the summary points of the executive summary on access
% request 2026-581, verbatim. Compile with pdflatex (twice).
\documentclass[10pt,letterpaper]{article}
\usepackage[T1]{fontenc}
\usepackage[utf8]{inputenc}
\usepackage[scaled=0.95]{helvet}
\renewcommand{\familydefault}{\sfdefault}
\usepackage[letterpaper,top=0.8in,bottom=0.8in,left=0.9in,right=0.9in]{geometry}
\usepackage{microtype}
\usepackage{enumitem}
\usepackage{titlesec}
\usepackage{fancyhdr}
\usepackage{xcolor}
\usepackage{booktabs}
\usepackage{longtable}
\usepackage{array}
\usepackage[hidelinks]{hyperref}
\urlstyle{same}
\newcommand{\docversion}{2026-10-16-16-30-27} % set by build.py (America/Toronto)
\hypersetup{pdfsubject={version \docversion}}

\setlength{\parindent}{0pt}
\setlength{\parskip}{0.5em}
\linespread{1.0}

\titleformat{\section}{\normalfont\large\bfseries}{\thesection}{0.6em}{}[\vspace{1pt}\titlerule]
\titlespacing*{\section}{0pt}{1.1em}{0.4em}
\titleformat{\subsection}{\normalfont\bfseries}{\thesubsection}{0.6em}{}
\titlespacing*{\subsection}{0pt}{0.8em}{0.25em}

\pagestyle{fancy}
\fancyhf{}
\renewcommand{\headrulewidth}{0pt}
\fancyfoot[L]{\footnotesize\color{gray}Access request 2026-581 -- summary points -- \pagesurl}
\fancyfoot[R]{\tiny\color{gray}v.\,\docversion}

% ---- Repository links: fill in once the documents are uploaded ----
\newcommand{\repourl}{https://github.com/OpenMobilityData/Terrebonne-access-requests}
\newcommand{\repo}[1]{\href{\repourl}{#1}}
% \doc{path-in-repo}{label}: link to a specific file in the repository
\newcommand{\doc}[2]{\href{\repourl/blob/main/#1}{#2}}
% GitHub Pages URL (Settings > Pages > deploy from main, root). Page anchors
% (#page=N) work there because Pages serves PDFs with the right content type;
% GitHub's blob viewer and raw.githubusercontent.com do not honour them.
\newcommand{\pagesurl}{https://openmobilitydata.github.io/Terrebonne-access-requests}
\expandafter\def\csname bundle01\endcsname{2026-581/raw/01_emails_mandate_counts.pdf}
\expandafter\def\csname bundle02\endcsname{2026-581/raw/02_framework_agreement_tender.pdf}
% \pg{bundle}{pages}: page reference, linked to the bundle at its first page
% \pgl{bundle}{anchor page}{displayed pages}: same with an explicit anchor
\newcommand{\pgl}[3]{\href{\pagesurl/\csname bundle#1\endcsname\#page=#2}{\textsf{\small[#1, p.~#3]}}}
\newcommand{\pg}[2]{\pgl{#1}{#2}{#2}}
\newcommand{\fillin}[1]{\textcolor{red}{[#1]}}
\newcommand{\quot}[1]{\begin{quote}\small #1\end{quote}}

\begin{document}

{\large\bfseries Access request 2026-581: what the borough's file shows}\hfill{\small October~9, 2026, rev.\ October~16}

\vspace{-0.2em}
{\small Summary points of the executive summary, verbatim. Full analysis, page references and the released documents:\\ \repo{\repourl}}
\vspace{0.4em}

\begin{enumerate}[nosep,itemsep=0.3em,leftmargin=1.4em]
  \item \begin{minipage}[t]{\linewidth}\textbf{The objective was fixed before the study.} On June~10, 2026 the
    borough's engineering division wrote to the City's street-design
    division that the elected officials wished to revisit the 2024 layout
    ``avec l'objectif de redonner la bande de stationnement qui avait \'et\'e
    retranch\'ee,'' that the borough's engineers had \emph{already} issued a
    technical opinion on the two concepts the councillors had sketched, and
    that the councillors wanted ``un avis officiel, neutre et externe.'' The
    mandate eventually given to EXP on August~20 is titled, in part,
    ``\'evaluer la place du stationnement sur rue, dans le but d'en augmenter
    le nombre.''\end{minipage}\vspace{0.35em}
\item \begin{minipage}[t]{\linewidth}\textbf{The City's own cycling unit warned about safety.} The
    Service de l'urbanisme et de la mobilit\'e returned an express analysis on
    June~17; the borough's own summary of it, written five days later, says
    the options ``pr\'esentent plusieurs compromis importants sur le plan de
    la s\'ecurit\'e, du confort des usagers vuln\'erables et de la coh\'erence
    du r\'eseau cyclable.'' That analysis is withheld. The fire department's
    opinion of July~13 is redacted in its entirety.\end{minipage}\vspace{0.35em}
\item \begin{minipage}[t]{\linewidth}\textbf{The winter measure is described as transitional; the
    mandate from the mayor's office was permanent.} On August~13 the borough
    mayor's chief of staff set the mandate in writing: a bidirectional lane
    on one side of Terrebonne, from Loyola to C\^ote-Saint-Luc, ``qui
    respecte les normes minimales pr\'evues par la loi,'' and ``le retour du
    stationnement d'un c\^ot\'e de la rue.'' No winter measure appears. The
    service request a week later added Volet~1, removal of cycling
    facilities in winter, as ``mesure transitoire pour l'hiver
    2026--2027,'' ahead of Volet~2, the permanent bidirectional lane. The
    winter conversion announced on October~6 is the first step of a
    year-round reduction of cycling space that has not been disclosed to
    the public.\end{minipage}\vspace{0.35em}
\item \begin{minipage}[t]{\linewidth}\textbf{The ``study'' is a six-week, \$14{,}755 technical note with a
    council date fixed in advance.} Service request August~20; offer
    August~25; purchase order August~26; kick-off September~3; draft
    September~17; final due October~2; announcement October~6. Borough staff
    told the consultant on August~18 that ``l'objectif de nos \'elus est de
    pouvoir agir sur vos conclusions pour un horizon d\'ebut octobre'' and on
    September~10 that the draft was needed by the 18th ``pour respecter notre
    calendrier du Conseil d'arrondissement du 5 octobre.''\end{minipage}\vspace{0.35em}
\item \begin{minipage}[t]{\linewidth}\textbf{The bicycle counts cited at council rest on the weakest
    data in the file and understate winter use; the file contained
    continuous counter data for last winter, documenting average winter usage
    at least five times higher than the cited figure. The more robust data
    were disregarded in the councillor's announcement.} The councillor told
    the meeting that ``we have counters everywhere, we have more counters
    than before. So we know exactly how many people use it, when and how,''
    and that ``we go from 700 during the good months to 8 \ldots\ So that's 4
    cyclists, right? Including me.'' The ``8,'' and the ``12'' he cited in
    answer to a question from the floor, come from the borough's February~2025
    count presentation: not bicycles per day or per hour, but bicycles
    counted using a temporary camera (Miovision) installed at the Cavendish
    and Grand intersections over six
    peak-period hours of a single day, January~29, 2025, in the lane's first
    winter, when clearing was often delayed, and \emph{a day on which Environment
    Canada had issued a winter weather advisory and residents were urged to
    postpone non-essential travel}; the same camera recorded cars on the
    Terrebonne approaches and pedestrians down by about a quarter from the
    November count, bicycles by 95\%. The slide states none of this. Before the first email in the file (June~10) and the
    mandate from the mayor's office (August~13), a continuous counter had
    already recorded every bicycle on the street through the whole of the
    following winter, a full year more recent; the borough's engineers
    offered its data to the City on June~10, and its table was in the
    consultant's hands in August: 4{,}891 passages over December--March
    (41 a day; 46 on weekdays) and a second summer at about 700 a day. For
    the question that bears on safety, exposure, the relevant figure is the
    total number of winter trips, not the count on one chosen day; that
    total was in the borough's possession and was not cited at council. The ``99\%'' drop cited at the meeting, 700
    against 8, sets the largest number on one slide, Grand in
    September~2024 counted on all four approaches, against the smallest on
    the next, Grand in January~2025 on the Terrebonne axis alone; compared
    like for like, the deck's own slides give 97--98\%, label January
    $-95$\%, and the counter gives about 93\% for the seasons entire. The
    two sources agree that the drop is steep;
    they differ in the absolute number, and only one of them yields ``four
    people.'' Whatever the
    explanation, a measure presented as data-driven rested on six hours of
    counting on a single day and ignored thousands of hours of more recent
    counting over an entire winter.\end{minipage}\vspace{0.35em}
\item \begin{minipage}[t]{\linewidth}\textbf{The need for additional winter parking was never measured.}
    Across the 508 pages there is no inventory of on-street spaces, no
    occupancy count, no survey of residents, and no estimate of the spaces
    the conversion would create; the ``50 to 60 spaces'' the councillor gave
    the APC-NDG at its consultation meeting, and repeated on CBC
    \emph{Daybreak} on October~9 (``at least 50 to 60 parking spots,''
    adding that ``it's a bit too early'' and that ``a few rounds during the
    workshops with the services'' remain before ``the final number''),
    appears nowhere in the release. The mandate asks the consultant to count
    the spaces an arrangement ``pourrait permettre de recr\'eer'' and to
    measure ``l'am\'elioration de la disponibilit\'e du stationnement''; the
    offer promises an analysis of ``l'impact sur le nombre de places'' and
    lists no field work; the consultant worked from aerial imagery. Whatever
    the note says about parking is a count of spaces that could be drawn,
    not a measurement of whether they are needed. Nothing in the file
    records how many winter nights parking near Terrebonne is actually
    full.\end{minipage}\vspace{0.35em}
\item \begin{minipage}[t]{\linewidth}\textbf{The case for removing a lane rests on a view that the lanes are
    wider than their bicycle traffic justifies; that view is the
    councillors' opinion, not the engineers' ruling.} The June~10 email
    reports the elected officials' concern that the lanes' width was ``jug\'ee
    disproportionn\'ee par rapport au volume de cyclistes observ\'e \ldots\
    selon les \'elus,'' and Councillor Teodoresco repeated the point to the
    APC-NDG at the consultation meeting, though not at the October~6 council
    meeting. No engineer in the file endorses it. The borough's own engineers
    describe the alternatives as ``une voie de circulation tr\`es \'etroite
    (moins de 4 m\`etres)'' and ``un am\'enagement bidirectionnel tr\`es
    \'etroit''; the City's cycling unit found ``plusieurs compromis
    importants''; and the August~13 mandate asks for a design ``qui respecte
    les normes minimales pr\'evues par la loi,'' which concedes that the June
    concepts did not. Lane width is set by the cyclist's operating envelope,
    by passing room and by the snow-clearing equipment, none of which depend
    on how many people ride; nothing in the file describes the existing
    lanes as wider than the standard requires.\end{minipage}\vspace{0.35em}
\item \begin{minipage}[t]{\linewidth}\textbf{A storm-only option was within the mandate.} EXP's offer
    lists among the scenarios to compare ``la fermeture en p\'eriode de
    d\'eneigement seulement.'' The councillor told the APC-NDG, at the meeting
    he later cited at council as his consultation with the Association, that
    storm-only parking had been ruled out because residents could not be
    informed of it; the option was not mentioned in his October~6 address,
    and the analysis on which the ruling-out rests is in the withheld note.\end{minipage}\vspace{0.35em}
\item \begin{minipage}[t]{\linewidth}\textbf{The exercise does not follow the City's own published method
    for reviewing a bike lane.} On February~10, 2026 the administration
    launched its audit of the cycling network on the principle ``recueillir
    les donn\'ees, analyser et corriger au besoin'': technical analysis, data
    collection and field observation through the peak season, led by the
    Service de l'urbanisme et de la mobilit\'e, with results expected at the
    end of 2026; the same release names Terrebonne among streets where the
    City is working with the borough on safety. On Terrebonne the objective
    was fixed before any analysis, the SUM's analysis was unfavourable and
    withheld, the mandate was to increase parking, no data were collected,
    and the change was announced while the City's audit was still under
    way.\end{minipage}\vspace{0.35em}
\item \begin{minipage}[t]{\linewidth}\textbf{The withholding of the technical report rests on a ground
    that expires with the decision.} The note and the City's analysis are withheld under
    section~39, which applies ``until a decision is made on the
    recommendation.'' On October~6 the councillor said the elected officials
    were ``now ready to give our orientation to the services'' and that the
    seasonal removal ``is what we're going to propose''; his signed
    statement adds that ``one thing is certain: residents will see changes
    before Christmas 2026.'' The \emph{Montreal Gazette} reported it the
    next day under the headline ``Terrebonne St.~will drop to one bike lane
    in winter to allow more parking, borough says'' and the subheading
    ``Councillor justifies decision by noting there's a sharp drop-off in
    cycling traffic during the winter''; the words are the newspaper's, not
    the councillor's. On October~9 he said on CBC \emph{Daybreak} that the
    final study ``still need[s] to be submitted to us'' and that the change
    goes to the borough council, to City council and back to the borough.
    Either the recommendation has been decided and
    the ground has lapsed, or it has not and the outcome was announced
    before the decision. A
    review may be sought from the Commission d'acc\`es \`a l'information
    within 30 days, that is by \textbf{November~8, 2026}.\end{minipage}\vspace{0.35em}
\end{enumerate}

\end{document}
